\emph{One-loop and scheme-level constraints on the selection datum.} Write \(X=\Phi^\dagger\Phi\) and \(r=c_2/c_3\) for the four couplings \((c_3,c_2,c_1,c_0)\) of the general \(U(4)\)-invariant potential \(c_3\operatorname{Tr}X^2+c_2(\operatorname{Tr}X)^2+c_1\operatorname{Tr}X+c_0\), so that Equation~\eqref{eq:V} is the slice \((\lambda/4,\,0,\,-\lambda v^2/2,\,\lambda v^4)\).
At tree level nothing fixes this slice: the matter term \(S_{\mathrm{matter}}\) of Equation~\eqref{eq:Smatrix} is a symbolic placeholder that has never been specified, and the minimal kinetic lift of Equation~\eqref{eq:Smatrix} to the bifundamental of Equation~\eqref{eq:Smicro} is invariant under \(\Phi\mapsto\Phi+\kappa\,\mathbb I\), a symmetry that forbids any potential, whereas Equation~\eqref{eq:V} is not shift-invariant; every local \(W(X)\) assigned in Presentation~B is therefore an added datum rather than an output of Equation~\eqref{eq:Smatrix}.
At one loop the scalar sector generates \((\operatorname{Tr}X)^2\): for the counterterm kernel \(\operatorname{Tr}(H_WH_W)\), with \(H_W\) the Hessian in the real components of \(\Phi\), one has \(\operatorname{Tr}(H_{\operatorname{Tr}X^2}^2)=256\operatorname{Tr}X^2+96(\operatorname{Tr}X)^2\) at \(n=4\) (ratio \(3/8\)), so \(\beta_{c_2}>0\) on the slice \(c_2=0\).
The gauge contribution is group data: over an orthonormal Hermitian \(u(N)\) basis the gauge loop gives exactly \(+N\operatorname{Tr}X^2\) with vanishing \((\operatorname{Tr}X)^2\) coefficient, while over an orthonormal \(su(N)\) basis it gives \((N-2/N)\operatorname{Tr}X^2+(1/N^2)(\operatorname{Tr}X)^2\), so the generated ratio of \((\operatorname{Tr}X)^2\) to \(\operatorname{Tr}X^2\) at \(N=4\) is \(1/56\) against \(0\) for \(U(4)\), and the group choice itself is not fixed by Equation~\eqref{eq:Smicro}.
With \(w=g^4/c_3^2\), \(v_4=\sum_f y_f^4/c_3^2\), and \(B_G\) the coefficient of the gauge contribution to \(\beta_{c_2}\) in units of \(w\) (\(B_G=0\) for \(u(4)\), \(1/16\) for \(su(4)\)), \(\beta_{c_3}=c_3^2(256+192r+4w-4v_4)/(32\pi^2)\), \(\beta_{c_2}=c_3^2(96+512r+640r^2+B_Gw)/(32\pi^2)\), and \(\beta_r=c_3P(r)/(32\pi^2)\) with \(P(r)=448r^2+(256+4(v_4-w))r+96+B_Gw\).
Renormalizable Yukawa couplings enter through \(\operatorname{Tr}X^2\) only (single cyclic-word identity), so \(P(0)=96+B_Gw>0\) is fermion-independent, the slice \(c_2=0\) is not a fixed ray and not an infrared attractor, and real fixed rays exist if and only if \((w-v_4-64)^2>10752+112B_Gw\): at \(B_G=v_4=0\) none occur for \(0\le w\le 64+16\sqrt{42}\approx167.7>16\pi^2\), and off-window at \(w=200\) the two exact rays are \(r=3/14\) and \(r=1\).
The \(c_1\) direction is scheme- and content-dependent: a sharp cutoff generates \(\delta c_1=(\Lambda^2/64\pi^2)(64c_3+136c_2+\ldots)\operatorname{Tr}X\) with the omitted gauge-ghost and Yukawa terms of content-dependent sign (bosonic loops positive, fermionic negative), Pauli--Villars replaces the quadratic sensitivity by a finite regulator-spectrum datum, and in a mass-independent minimal-subtraction scheme \(\beta_{c_1}=c_1\gamma\) carries no additive term, so \(c_1=0\) is an exact renormalization surface rather than a derived value.
The constant-\(\Theta\) Moyal structure generates no potential either: every \(p\)-th star term carries at least one derivative on each factor, so at zero gradient the star potential collapses to the ordinary potential to all orders, the \(O(\Theta)\) correction being a total derivative and the \(O(\Theta^2)\) correction a pure derivative operator.
Scope: within this hypothesis class (tree-level Presentation~A with unspecified matter term or minimal shift-symmetric lift, a local Presentation-B potential \(W(X)\) for \(n=4\) bifundamental scalars, renormalizable gauge and Yukawa couplings, group choice as input, fixed-\(\Theta\) deformation, one-loop accuracy) these computations constrain but do not select the slice of Equation~\eqref{eq:V}; closing the selection would require an explicit \(S_{\mathrm{matter}}\) or lift, or a specified group, fermion content and ultraviolet boundary condition on \((c_3,c_2,c_1)\) at the cutoff \(\LUV\) of Equation~\eqref{eq:LUV-def}, or a non-renormalizable double-trace completion, none of which Equation~\eqref{eq:Smatrix} supplies; a dynamical-\(\Theta\) or strongly coupled branch is not evaluated here, and none of the listed data, if specified, is excluded from realizing the slice.
The selection of Equation~\eqref{eq:V} by Presentation~A therefore remains an input of the matching discussed in Section~\ref{sec:ch3-relate}: it is constrained but underdetermined as datum, and nothing in this paragraph closes Problem~D, Problem~D1, or Problem~A.
